\section{Introduction}
\label{sec:introduction}

\subsection{Motivation}
\label{sec:motivation}

Training humanoid whole-body control policies via imitation learning
requires large volumes of robot-executable reference motion. Since motion
capture is collected on humans, every such pipeline needs a retargeting step
that maps human kinematics onto a specific robot's morphology while
respecting its joint limits, self-collision geometry, and physical
plausibility. General Motion Retargeting (GMR) has emerged as a widely used
tool for this step: a per-frame inverse-kinematics solver (QP-based, built on
\texttt{mink}) that tracks human landmark targets on the robot subject to
hard joint limits.

GMR's own published evaluation, however, explicitly excludes motions with
complex floor contact: the authors state that they ``do not include motions
with complex interaction with the environment, such as crawling or getting
up from the floor''~\cite{araujo2025gmr}. This is a deliberate, documented
scoping choice, not an oversight. But it leaves an open question with direct
practical consequences: humanoid robots that fall, or that need to
manipulate objects from the ground, or that recover from being knocked over,
need exactly this class of motion as training data. If the standard
retargeting tool cannot handle it, either the training pipeline silently
drops this data, or it trains on badly-retargeted data without knowing it.

This is not a narrow gap. A broader survey of the human-to-humanoid
retargeting literature (contact-aware sequence retargeters, kinodynamic
refiners, learning-based retargeters --- Section~\ref{sec:related-work})
finds the same exclusion pattern repeated: contact modeling in this line of
work stops at feet and hands, target-side feasibility editing addresses foot
placement and root trajectories but not knee, pelvis, or forearm support,
and every hardware demonstration of a real floor transition to date treats
it as an \textit{emergent behavior of a learned policy} rather than the
output of a retargeter. We share this observation with a larger,
forward-looking research plan for this line of work (\texttt{paperIdea3.md},
in this repository) that scopes a full contact-template, hardware-validated
system as a 2027 target. The present paper is deliberately narrower: it
validates the first, load-bearing claim that plan depends on --- that GMR's
exclusion is a real, measurable gap, and that a targeted contact-aware layer
on top of its own solver closes most of it kinematically --- before
committing to the larger system. We return to that larger scope in
Section~\ref{sec:further-future-work}.

\subsection{Related Work}
\label{sec:related-work}

\noindent\textbf{Per-frame kinematic and teleoperation retargeters.}
\cite{darvish2019geometric} established geometric per-frame IK retargeting
validated across multiple humanoid platforms and remains a widely
referenced baseline pattern in this literature. Such systems enforce joint
limits and simple center-of-mass or zero-moment-point constraints per frame,
but have no explicit multi-link floor-contact modeling and no whole-clip
temporal optimization --- smoothness is whatever the per-frame solve and a
downstream controller happen to produce.

\noindent\textbf{Offline whole-sequence optimization retargeters.}
\cite{ayusawa2017morphing} solve human-model identification,
human-to-robot morphing, and robot motion planning as one simultaneous
trajectory optimization, enforcing joint limits and a simplified
linear-inverted-pendulum center-of-mass/zero-moment-point balance
constraint. This is the same style of quasi-static balance check an earlier
phase of this project built and deliberately left disabled by default
(Section~\ref{sec:limitations}), because a kinematic center-of-mass check
cannot distinguish genuine imbalance from momentum-carried dynamic posture
--- a limitation of this entire class of check, not particular to our own
attempt at it. \cite{otani2017adaptive}, building on
\cite{difava2016multicontact}, extend this line to explicit multi-contact
retargeting (hands, feet, and environment contacts jointly) and remain a
widely-cited foundation for later contact-aware pipelines. All three are,
like ours, target-side geometric or optimization-based methods rather than
RL- or controller-side ones; none evaluate on floor-contact postural
transitions (kneeling, prone/supine, crawling, getting up) as a first-class
case, and each assumes the robot stays continuously balanced and in known
contact with the ground --- an assumption every clip in our floor-contact
class violates by construction.

\noindent\textbf{Modern contact-aware sequence retargeters.}
\cite{jeong2025rigunification} unifies diverse motion sources onto a
canonical rig, then refines trajectories to enforce foot-contact constraints
and stability, validated on 12 simulated and 3 real humanoids --- the most
recent, and methodologically closest, prior work to ours. It is explicitly
scoped to foot-fixed, upper-body-expressive motion: its own feasibility
constraint requires both feet to remain in continuous ground contact
throughout, which structurally excludes exactly the motion class this paper
targets, where the robot must lift, reposition, and re-plant feet, or bear
weight through knees, pelvis, or a prone torso. \cite{yang2025omniretarget}
(OmniRetarget) preserves human-environment-object contact relationships
through an interaction mesh and Laplacian deformation, and is highly
influential as a data-generation backbone for later reinforcement-learning
pipelines; its contact modeling, like ours, is target-side and geometric,
but its demonstrated scope is loco-manipulation and terrain interaction, not
floor-contact postural transitions. \cite{jeong2025core} (CoRe) is
philosophically closest to our own position --- fix the geometry before
anything downstream has to compensate for it --- applying contact-aware
optimization refinement before an RL policy, but it too is scoped to
loco-manipulation contact (foot sliding, floating), not knee, pelvis, or
prone floor support.

\noindent\textbf{Retargeting as a reinforcement-learning data-generation
front end.} A large and fast-growing body of recent work --- TeleGate
\cite{li2026telegate}, the Humanoid Manipulation Interface
\cite{nai2026humi}, HumanPlus \cite{fu2024humanplus}, and the tracking-policy
ecosystem those and GMR feed, including BeyondMimic
\cite{liao2025beyondmimic} --- demonstrates that floor-contact transitions
(standing up, fall recovery, kneeling) are achievable on real hardware
today. In every one of these systems the transition is a \textit{learned
policy behavior}, executed from a reference trajectory the retargeter did
not need to get right at the contact level; the retargeter's job is only to
produce a plausible-enough reference for the policy to correct. Our
contribution is upstream of and complementary to this line of work: we ask
whether the retargeted reference itself can be made contact-correct by
construction, kinematically, before any RL policy sees it, matching this
literature's own standard practice of validating reference-motion quality
as a step distinct from policy training (Section~\ref{sec:scope}) --- and
directly motivating the policy-training evaluation planned for this paper
(Section~\ref{sec:status-next-steps}).

\noindent\textbf{A kinematic-only correction is a deliberate scope choice,
not an oversight.} An alternative to our approach is to add dynamics
directly to the correction layer: DynaRetarget
\cite{dhedin2026dynaretarget} uses sampling-based trajectory optimization to
convert imperfect kinematic retargets into dynamically feasible motions for
loco-manipulation. We do not take that route here, so that the correction
layer stays solver-agnostic and inexpensive enough to run offline over a
full 77-clip corpus without kinodynamic optimization or a physics simulator
in the loop (Section~\ref{sec:overview}). Whether the resulting kinematic,
geometric contact-correctness guarantee is also sufficient for dynamic
trackability is exactly the open question the planned policy-training
evaluation (Section~\ref{sec:status-next-steps}) answers --- this paper's
contribution stops at the kinematic guarantee, deliberately.

\noindent\textbf{Where this leaves the gap.} Across all four groups above,
floor-contact postural transitions are either out of scope by explicit
assumption (continuous ground contact, feet fixed), addressed only for
feet, hands, or manipulated objects rather than knees, pelvis, or a prone
torso, or resolved at the policy level rather than the retargeter level. We
are not aware of prior work that treats crawling, kneeling, prone/supine
lying, or getting up as a first-class kinematic retargeting target with an
explicit, non-gameable contact-correctness guarantee
(Section~\ref{sec:grounding}). That is the gap this paper closes.

\subsection{The Gap}
\label{sec:gap}

We first quantify the gap. Running GMR out-of-box on LAFAN1 clips it did not
validate against --- clips involving crawling, falling, and getting up ---
we measure 13--16~cm of floor penetration on up to 91\% of frames, an
order-of-magnitude worse than GMR's own clean-locomotion clips (1~cm, 0.3\%
of frames). This is confirmed at full-corpus scale (77 LAFAN1 clips, split
34 floor-contact / 43 locomotion by a multi-surface human contact detector,
not just a hip-height heuristic --- hip-height alone undercounts floor
contact by missing hand/knee-supported poses during brief stumbles).

We also show that GMR's own published mitigation --- a single clip-wide
constant height offset --- cannot fix this: applying it removes most of the
penetration but leaves the robot floating 6--18~cm above the floor on
average during exactly the frames that should show tight ground contact,
because a single rigid shift cannot simultaneously satisfy a clip's standing
phase and its lying phase. We built a specific metric to make this failure
mode visible and non-gameable (Section~\ref{sec:grounding}), since floor
penetration alone and floating alone can each be minimized independently by
a shift in the wrong direction.

\subsection{Contributions}
\label{sec:contributions}

\begin{enumerate}
  \item \textbf{A contact-aware correction layer}, in two instantiations,
    that sits entirely on top of GMR's own solver with zero changes to
    GMR's core tracking logic: a multi-stage per-frame pipeline (floor and
    self-collision correction, held-aware temporal smoothing, rate-limited
    correction) and a global whole-trajectory pipeline (a single sparse QP
    jointly enforcing floor/collision correctness, contact anchoring, and
    temporal smoothness), both terminating in a local grounding envelope
    with an algebraic zero-penetration guarantee.
  \item \textbf{A reference-free physics-plausibility evaluation suite} ---
    no motion-capture ground truth on the robot is available, so every
    metric we report (floor penetration, self-collision, contact
    correctness, tracking fidelity, smoothness) is computed from the
    retargeted motion and the robot's own geometry alone.
  \item \textbf{A non-gameable joint contact-correctness metric}
    (Section~\ref{sec:grounding}): a naive constant Z-shift can
    independently minimize floor penetration or eliminate floating, but not
    both at the same held frame simultaneously --- our metric requires
    both, closing a gaming loophole present in simpler, single-axis
    versions of the same check.
  \item \textbf{Full 77-clip LAFAN1 corpus validation}, floor/locomotion
    class split, against GMR's own published height-fix baseline.
  \item \textbf{An honest architectural finding and open trade-off}:
    per-frame independent contact correction can select different,
    equally-valid null-space solutions on adjacent frames, producing a
    visible ``branch-flip'' artifact that aggregate corpus metrics do not
    surface --- caught only by watching renders. We trace this to the
    correction's lack of temporal coupling, characterize a partial
    per-frame mitigation, and report results from an alternative
    whole-trajectory optimization formulation that removes the artifact
    structurally but currently trails the per-frame approach on contact
    precision. We present this as an open problem, not a solved one.
\end{enumerate}

\subsection{Scope}
\label{sec:scope}

This is a module/polish paper, not a new solver. We do not replace GMR's own
tracking IK; we validate that a targeted contact-correctness layer on top of
it closes a real, documented gap in its applicability, and we report the
result honestly including where it does not yet fully close. Kinematic
plausibility --- physically executable, contact-correct, smooth motion ---
is necessary but not sufficient for downstream use, consistent with the
standard practice in this literature of validating retargeting quality
before the separate, more expensive step of training a tracking policy on
it. Imitation-policy training on the retargeted corpus, which closes that
gap, is in progress at the time of writing and its results are planned as
part of this paper (Section~\ref{sec:status-next-steps}).
